Let \(X_i=x_0\) almost surely and let \((G_i,S_i)\) be uniform on
\[
 \{(2,0),(2,1),(3,0),(3,1),(\infty,0),(\infty,1)\}.
\]
On a common probability space, take \(U_i\) and \(E_i\) to be independent Rademacher variables. Take \(H_i\), independent of \((U_i,E_i,G_i,S_i)\), with
\[
 P(H_i=1)=\frac{1+\rho}{2},
 \qquad
 P(H_i=-1)=\frac{1-\rho}{2}.
\]
Also take \((U_i,E_i,H_i)\) independent of \((G_i,S_i)\). Define
\[
 V_i=
 \begin{cases}
 E_i,&G_i\in\{2,3\},\\
 U_iH_i,&G_i=\infty.
 \end{cases}
\]
These independence statements imply, in every cohort-eligibility cell,
\[
 E[U_i\mid G_i=a,S_i=s]=E[V_i\mid G_i=a,S_i=s]=0,
 \qquad
 E[U_i^2\mid G_i=a,S_i=s]=E[V_i^2\mid G_i=a,S_i=s]=1,
\]
and
\[
 E[U_iV_i\mid G_i=a,S_i=s]
 =\begin{cases}
 0,&a\in\{2,3\},\\
 \rho,&a=\infty.
 \end{cases}
\]
Indeed \(U_i\) and \(E_i\) are independent and centered, whereas \(U_iV_i=H_i\) in the never-enabled cohort.

Put \(r=1/8\), \(Z_{i1}=0\), \(Z_{i2}=rU_i\), and \(Z_{i3}=rV_i\). For \(s\in\{0,1\}\) and \(t\in\{1,2,3\}\), define
\[
 Y_{it}(\infty,s)=Z_{it},
\]
and, for \(a\in\{2,3\}\), define
\[
 Y_{it}(a,s)=Z_{it}+s\delta\mathbf 1\{t\ge a\}.
\]
Finally set \(Y_{it}=Y_{it}(G_i,S_i)\). Consistency holds by construction. For eligible units, the added term is zero before adoption, proving no anticipation. For ineligible units it is identically zero for every enabling date, proving no spillover to ineligibles.

Under never enabling, both eligibility states have the outcome history \((Z_{i1},Z_{i2},Z_{i3})\). Conditional centering therefore makes both sides of every retained pairwise conditional DDD trend restriction zero. The two target effects for cohort \(2\) and all retained conditional DDD contrasts equal \(\delta\):
\[
 d_{(2,2),3}=d_{(2,2),\infty}
 =d_{(2,3),\infty}=\delta,
 \qquad
 \theta_{(2,2)}=\theta_{(2,3)}=\delta.
\]
Consequently every target-cell contrast term \(d_{k,c}(X_i)-\Theta_k(P)\) is zero. The post-adoption effect is nonzero because \(\delta\ne0\), but it disappears upon residual centering, so it does not affect the score covariance.

All six cohort-eligibility probabilities equal \(1/6>1/8\), so overlap holds. Moreover,
\[
 \max_t|Y_{it}|\le r+|\delta|<\frac14,
\]
which gives the outcome fourth-moment bound with \(C_4=1\).

Because \(p_{2,1}=\pi_{2,1}=1/6\), every nonzero residual coefficient in a retained score equals six times its DDD sign. Write
\[
 j_1=((2,2),3),\qquad
 j_2=((2,2),\infty),\qquad
 j_3=((2,3),\infty).
\]
The residual for \(j_1\) is \(rU_i\) in the four cells with cohort \(2\) or \(3\); the residual for \(j_2\) is \(rU_i\) in the four cells with cohort \(2\) or \(\infty\); and the residual for \(j_3\) is \(rV_i\) in the four cells with cohort \(2\) or \(\infty\). Mutually exclusive cells have zero covariance. A shared cell contributes
\[
 \frac16\,6^2r^2\operatorname{Cov}(U_i,V_i\mid G_i=a,S_i=s)
 =6r^2\operatorname{Cov}(U_i,V_i\mid G_i=a,S_i=s)
\]
when the residual signs agree. There are four cells on each diagonal, two cohort-\(2\) cells shared by \(j_1,j_2\), no nonzero-covariance cell shared by \(j_1,j_3\), and two never-enabled cells with covariance \(\rho\) shared by \(j_2,j_3\). Since \(r^2=1/64\), Lemma \(\ref{lem:shared-covariance-expansion}\) gives
\[
 \Omega(\rho)=r^2
 \begin{pmatrix}
 24&12&0\\
 12&24&12\rho\\
 0&12\rho&24
 \end{pmatrix}
 =
 \begin{pmatrix}
 3/8&3/16&0\\
 3/16&3/8&3\rho/16\\
 0&3\rho/16&3/8
 \end{pmatrix}.
\]
Its eigenvalues are
\[
 \frac3{16}\left(2,\ 2-\sqrt{1+\rho^2},\ 2+\sqrt{1+\rho^2}\right).
\]
The restriction \(|\rho|<\sqrt7/3\) gives \(\sqrt{1+\rho^2}<4/3\). Hence the smallest eigenvalue is strictly greater than \(1/8\), while the largest is less than \(5/8<8\). Thus
\[
 \frac18I_3\preceq\Omega(\rho)\preceq8I_3.
\]
Every defining condition of \(\mathcal M\) is now verified, so \(Q_{\rho,\delta}\in\mathcal M\).

The first covariance block has equal diagonal entries and symmetric off-diagonal entries. Its componentwise optimum therefore assigns equal weights, while the second block is a singleton:
\[
 b^{\mathrm{sep}}=
 \begin{pmatrix}\omega_1/2\\\omega_1/2\\\omega_2\end{pmatrix}.
\]
For \(N=(1,-1,0)^{\prime}\), direct multiplication gives
\[
 N^{\prime}\Omega N=\frac38,
 \qquad
 N^{\prime}\Omega b^{\mathrm{sep}}=-\frac{3\rho\omega_2}{16}.
\]
The contrast correction from Theorem \(\ref{thm:separate-joint-frontier}\) is therefore
\[
 -(N^{\prime}\Omega N)^{-1}N^{\prime}\Omega b^{\mathrm{sep}}
 =\frac{\rho\omega_2}{2}.
\]
Thus
\[
 b^{\mathrm{joint}}
 =b^{\mathrm{sep}}+\frac{\rho\omega_2}{2}N
 =\begin{pmatrix}
 \omega_1/2+\rho\omega_2/2\\
 \omega_1/2-\rho\omega_2/2\\
 \omega_2
 \end{pmatrix},
\]
and
\[
 V_{\mathrm{sep}}-V_{\mathrm{joint}}
 =\frac{(3\rho\omega_2/16)^2}{3/8}
 =\frac{3\rho^2\omega_2^2}{32}>0.
\]
Both \(\rho\) and \(\delta\) range over nonempty open sets, and all inequalities above are strict throughout those ranges. This proves the asserted nondegenerate open family.